% 英文摘要

Ultra-high vacuum (UHV) is a fundamental prerequisite for modern high-end scientific experiments and advanced manufacturing technologies, and its measurement accuracy directly determines the reliability of related experimental results and process performance. Conventional ionization-based vacuum gauges rely on gas ionization processes and are susceptible to gas composition, surface effects, and device aging, leading to inherent systematic errors and limited long-term stability. These limitations hinder their ability to meet the requirements of high-precision and traceable vacuum metrology. 

Quantum vacuum sensing based on cold-atom collision-induced loss utilizes the interactions between background gas particles and trapped cold atoms, establishing a direct relationship between the atomic loss rate and the vacuum pressure. The measurement is inherently traceable to fundamental physical constants, offering high accuracy and excellent long-term stability. In recent years, this approach has emerged as a promising direction in vacuum metrology. However, existing studies have primarily focused on proof-of-principle demonstrations and metrological model development, while challenges in system miniaturization, engineering integration, and practical deployment remain significant, limiting its broader application.

In this work, we investigate the design and experimental validation of a miniaturized cold-atom vacuum sensing system. To address the large footprint and structural complexity of conventional cold-atom apparatus, an integrated system design is developed, including a compact optical setup, a hybrid permanent-magnet and electromagnetic coil magnetic field configuration, and a miniaturized ultra-high vacuum chamber. Based on a rubidium atomic platform, a complete cold-atom experimental system is constructed and magneto-optical trap (MOT) loading is achieved. Building upon this, lifetime measurements of atoms in a magnetic trap are carried out by monitoring the temporal decay of atom number, from which the loss rate is extracted and quantitatively related to the background gas pressure. Furthermore, the consistency, repeatability, and dominant sources of systematic uncertainty are systematically analyzed.

Experimental results demonstrate that the proposed miniaturized cold-atom vacuum sensing system can reliably realize MOT loading and exhibits good measurement consistency and repeatability. This work provides experimental evidence and technical guidance for the miniaturization of cold-atom-based vacuum sensors, and contributes to the advancement of quantum vacuum metrology toward practical and engineering applications.
